\section{第~\ref{sec:pain}~章　疼痛}

痛觉传感器、两根导线、脊髓闸门、下行的音量旋钮、敏化和对注意力的夺取都已直接测量；闸门和敏化的公式是本书的简化模型；机器选择警报音量的规则是假说。

{\footnotesize
\setlength{\tabcolsep}{3pt}\renewcommand{\arraystretch}{1.15}
\begin{longtable}{>{\raggedright}p{0.5cm}>{\raggedright}p{4.0cm}>{\raggedright}p{5.6cm}>{\raggedright}p{2.3cm}>{\raggedright\arraybackslash}p{1.7cm}}
\toprule
序号 & 命题 & 实验及测量内容 & 文献 & 状态 \\
\midrule\endfirsthead
\toprule
序号 & 命题 & 实验及测量内容 & 文献 & 状态 \\
\midrule\endhead
1 & 高阈值：不同痛觉传感器的热阈值从$41^\circ$到$46$--$48^\circ$ & 单纤维记录。猴的皮肤：慢（C）传感器的中位阈值为$41^\circ$，第二类快传感器为$46^\circ$，第一类高于$53^\circ$。人的皮肤：同时响应压力的C传感器为$40.7^\circ$，不响应压力的为$48^\circ$。 & \cite{Treede1995heat, Weidner1999cfib} & 已测量 \\
2 & 痛觉传感器的适应远弱于触觉传感器，轻度损伤后变得更敏感；强烈加热后的减弱只在一类中测到 & 皮肤轻度烫伤（$50^\circ$，$100\,\text{s}$）：$5$--$10$\,min后，人的痛阈和猴C传感器的阈值降低$2$--$6^\circ$；其他皮肤传感器没有这种偏移。第二类快传感器在强烈加热后响应受到抑制。与触觉传感器适应的比较是一般性估计，此处没有单独的实验支持。 & \cite{LaMotte1982hyper, Treede1995heat} & 已测量 \\
3 & 两根导线：快的$5$--$30$\,m/s，慢的$0.5$--$2$\,m/s；到达时间$t=L/v$~\eqref{eq:paintime} & 传导速度：猴的快痛觉传感器平均$15$和$25$\,m/s（两类）；人的C传感器$0.8$--$1.0$\,m/s。$L=1$\,m时，分别为几十毫秒和约一秒。 & \cite{Treede1995heat, Weidner1999cfib} & 已测量 \\
4 & 疼痛来两次：先快而准，后钝而久 & 人前臂受热的反应时间：随温度升高从$1100$降到$400\ms$；强烈加热有毛皮肤产生双重疼痛，手掌则不会。第一痛只能由快于$6$\,m/s的纤维传送——按潜伏期与之对应的是第二类快传感器，而手掌中没有这类传感器。角色分工（“哪里”和“有多糟”）是一种解读。 & \cite{CampbellLaMotte1983first, Treede1995heat} & 已测量 \\
5 & 闸门：脊髓输出神经元接收疼痛和触摸，抑制性中间神经元被触摸兴奋（图~\ref{fig:gate}） & 闸门方案作为理论提出。小鼠，遗传标记并去除神经元群：表达生长抑素的兴奋性神经元是机械痛所必需的；表达强啡肽的抑制性神经元阻止来自触觉传感器的输入到达它们——没有这些神经元，轻触就会引起疼痛。 & \cite{MelzackWall1965, Duan2014} & 已测量 \\
6 & 触摸减轻疼痛 & 人，手臂上的激光痛觉脉冲：同时触摸降低疼痛评分和分辨脉冲能量的能力；在同一皮节内，触摸点与疼痛点距离越远，效应越弱。 & \cite{Mancini2014touch} & 已测量 \\
7 & 闸门理论的假设：剧痛自己关闭抑制性中间神经元，闸门大开 & 本章标注为原始理论的假设。小鼠研究描述了闸门如何阻止触觉进入疼痛通道；其中没有显示疼痛关闭中间神经元。 & \cite{MelzackWall1965} & 假说 \\
8 & 闸门~\eqref{eq:gate}：无触摸阈值$0.18$，有触摸$0.86$，$g=2$时$0.11$；$\nu=1$时触摸使输出从$1$降到$0.25$，$\nu=2$时无效；触摸本身不通过（图~\ref{fig:gatecurves}） & 用正文中的权重按\texttt{models/\allowbreak pain\_\allowbreak gate.py}计算，再现了所有这些数字。 & \texttt{models/\allowbreak pain\_\allowbreak gate.py} & 模型 \\
9 & 来自脑干的下行通路双向改变闸门增益 & 大鼠延髓，加热尾部时记录：一些神经元在缩尾之前放电（“开”细胞），另一些沉默（“关”细胞）；微弱刺激该区域可抑制反射。综述：同样的回路既易化也抑制疼痛传递，阿片类经由它们起作用。 & \cite{Fields1983rvm, Fields2004} & 已测量 \\
10 & 对缓解的预期经阿片通道减轻疼痛 & 急性疼痛患者：对那些因预期缓解而疼痛减轻的人，阻断阿片受体使疼痛回到其他人的水平；对其他人，阻断不改变疼痛。 & \cite{Levine1978} & 已测量 \\
11 & 大脑根据中断的收益选择警报音量 & 没有测量过音量选择规则的实验。 & --- & 假说 \\
12 & 敏化：损伤后闸门输出增大，阈值降低，部分来自脊髓；在损伤性刺激停止后仍然持续 & 大鼠：皮肤损伤后屈曲反射阈值下降，响应增大；电生理分析显示，部分增长产生于脊髓本身。损伤性刺激引起持久的感觉改变，比刺激本身持续得更久。本章未给出确切的消退时间。 & \cite{Woolf1983} & 已测量 \\
13 & 敏化~\eqref{eq:sensitization}是正反馈；环路强时，警报可能在愈合后自锁 & 来自模型~\eqref{eq:sensitization}（本章末的习题）。没有实验表明愈合后的持续疼痛正是这种形式的自锁环路。 & 本章推导 & 假说 \\
14 & 疼痛是负奖励，从疼痛中学习遵循时间差分误差 & 人，断层扫描，疼痛先兆序列：腹侧纹状体和前岛叶皮层的活动与时间差分学习信号吻合。猴：一部分\nt{DOPA}神经元被不愉快的吹气及其先兆兴奋，另一部分被抑制。 & \cite{Seymour2004, Matsumoto2009} & 已测量 \\
15 & 疼痛中断当前工作 & 人，电痛刺激加声音探测：疼痛开始后$100\ms$时对探测的反应明显变慢，$1.5\,\text{s}$时已与对照无差别。综述把这类实验归纳为注意力夺取模型。 & \cite{Crombez1996disrupt, EcclestonCrombez1999} & 已测量 \\
16 & 缩回反射在脊髓中闭合 & 大鼠：背角深层神经元重复了单块肌肉缩回反射的空间模式；无法从颈段逆向激活它们，即它们是脊髓中间神经元。 & \cite{Schouenborg1995wr} & 已测量 \\
\bottomrule
\end{longtable}}
